\begin{thebibliography}{9999}
	
	% --- 在这里插入间距调整命令 ---
	\setlength{\itemsep}{0.1em} % 调整每条文献之间的距离（数值越大，间距越宽）
	\setlength{\parsep}{0pt}    % 调整同一条文献内部不同段落的间距（通常设为0pt即可）
	% ------------------------------

	\bibitem{Gelfond} Gelfond, A., Sur les nombres transcendants, C. R. Acad. Sci. Paris, 189, p. 1224-1226 (1929); Sur le septième problème de Hilbert, Bull. Acad. Sci. U.R.S.S. (7), p. 625-630 (1930).

	\bibitem{Hermite} Hermite, Ch., Sur la fonction exponentielle, Oeuvres III, p. 150-181.

	\bibitem{Hilbert} Hilbert, D., Ueber die Transcendenz der Zahlen $e$ und $\pi$, Math. Ann. 45, p. 216-219 (1895).

	\bibitem{Hille} Hille, E., Gelfond's solution of Hilbert's seventh problem, Amer. Math. Monthly 49, p. 654-661 (1942).

	\bibitem{Hurwitz} Hurwitz, A., Beweis der Transcendenz der Zahl $e$, Math. Ann. 45, p. 220-222 (1895).

	\bibitem{Koksma} Koksma, J. F., Diophantische Approximationen, Ergebnisse der Mathematik und ihrer Grenzgebiete 4 (1936).

	\bibitem{Kuzmin} Kuzmin, R. O., Ob odnom novom klasse transcendentnyh cisel, Bull. Acad. Sci. U.R.S.S. (7), p. 585-597 (1930).

	\bibitem{Lambert} Lambert, J. H., Mémoire sur quelques propriétés remarquables des quantités transcendantes circulaires et logarithmiques, Opera mathematica II, p. 112-159.

	\bibitem{Legendre} Legendre, A. M., Éléments de géométrie, Note IV.

	\bibitem{Lindemann} Lindemann, F., Ueber die Zahl $\pi$, Math. Ann. 20, p. 215-225 (1882).

	\bibitem{Loewy} Loewy, A., Über Matrizen und Differentialkomplexe, Math. Ann. 78, p. 1-51 (1918).

	\bibitem{Maier} Maier, W., Potenzreihen irrationalen Grenzwertes, J. reine angew. Math. 156, p. 93-148 (1927).

	\bibitem{Schneider} Schneider, Th., Transzendenzuntersuchungen periodischer Funktionen I., J. reine angew. Math. 172, p. 65-69 (1935); Arithmetische Untersuchungen elliptischer Integrale, Math. Ann. 113, p. 1-13 (1937); Zur Theorie der Abelschen Funktionen und Integrale, J. reine angew. Math. 183, p. 110-128 (1941).

	\bibitem{Siegel} Siegel, C. L., Über einige Anwendungen diophantischer Approximationen, Abh. Preuss. Akad. der Wissensch., Phys. math. Kl., Jahrg. 1929, Nr. 1, 70p.

	\bibitem{Stridsberg} Stridsberg, E., Sur quelques propriétés arithmétiques de certaines fonctions transcendantes, Acta math. 33, p. 233-292 (1910).

	\bibitem{Weierstrass} Weierstrass, K., Zu Lindemann's Abhandlung, "Über die Ludolph'sche Zahl", Mathematische Werke II, p. 341-362.

\end{thebibliography}